% =====================================================================
%  Study Notes for MTP viva
%  Author: Sakhare Yash Balram (2022CH71496)
%
%  This is a self-study reference, NOT the main report.  It is meant to
%  be a single document that takes you from "what is OpenFOAM" through
%  to "exactly what is in every file of my case", with all the numbers
%  and code that you should be able to talk through in the viva.
% =====================================================================

\documentclass[11pt,a4paper]{article}

\usepackage[T1]{fontenc}
\usepackage{lmodern}
\usepackage[utf8]{inputenc}
\usepackage{microtype}
\usepackage[margin=22mm]{geometry}
\usepackage{setspace}
\onehalfspacing

\usepackage{amsmath,amssymb}
\usepackage{siunitx}
\usepackage{graphicx}
\usepackage{float}
\usepackage{booktabs}
\usepackage{xcolor}
\usepackage{enumitem}
\usepackage[explicit]{titlesec}
\titleformat{\section}{\sffamily\Large\bfseries}{\thesection.}{0.5em}{#1}
\titleformat{\subsection}{\sffamily\large\bfseries}{\thesubsection.}{0.5em}{#1}
\titleformat{\subsubsection}{\sffamily\normalsize\bfseries}{\thesubsubsection.}{0.5em}{#1}
\titleformat{\paragraph}[runin]{\sffamily\normalsize\bfseries}{}{0pt}{#1}

\usepackage{listings}
\definecolor{codebg}{rgb}{0.97,0.97,0.97}
\definecolor{codekey}{rgb}{0.10,0.30,0.65}
\definecolor{codestr}{rgb}{0.55,0.20,0.10}
\definecolor{codecmt}{rgb}{0.40,0.50,0.40}
\lstset{
    backgroundcolor=\color{codebg},
    basicstyle=\ttfamily\footnotesize,
    keywordstyle=\color{codekey}\bfseries,
    commentstyle=\itshape\color{codecmt},
    stringstyle=\color{codestr},
    showstringspaces=false,
    numbers=left,
    numberstyle=\tiny\color{gray},
    stepnumber=1,
    breaklines=true,
    breakatwhitespace=true,
    frame=single,
    framesep=4pt,
    rulecolor=\color{black!20},
    captionpos=b,
    xleftmargin=10pt,
}
\lstdefinelanguage{foam}{
    morekeywords={FoamFile,version,format,class,object,location,
        application,startFrom,startTime,stopAt,endTime,deltaT,writeControl,
        writeInterval,purgeWrite,writeFormat,writePrecision,writeCompression,
        timeFormat,timePrecision,runTimeModifiable,adjustTimeStep,maxCo,
        maxDeltaT,functions,libs,fields,operation,regionType,name,writeFields,
        log,type,internalField,uniform,boundaryField,wall,sym,main_inlet,
        outlet,branch_inlet,fixedValue,zeroGradient,inletOutlet,inletValue,
        codedFixedValue,symmetryPlane,kqRWallFunction,omegaWallFunction,
        nutkWallFunction,fixedFluxPressure,prghPressure,
        pressureInletOutletVelocity,turbulentIntensityKineticEnergyInlet,
        turbulentMixingLengthFrequencyInlet,intensity,mixingLength,value,
        dimensions,solvers,solver,smoother,preconditioner,tolerance,relTol,
        nSweeps,maxIter,GAMG,PCG,DIC,DILU,PBiCGStab,smoothSolver,
        symGaussSeidel,GaussSeidel,diagonal,PIMPLE,SIMPLE,
        momentumPredictor,nOuterCorrectors,nCorrectors,
        nNonOrthogonalCorrectors,relaxationFactors,equations,
        ddtSchemes,gradSchemes,divSchemes,laplacianSchemes,
        interpolationSchemes,snGradSchemes,fluxRequired,
        Euler,steadyState,backward,CrankNicolson,Gauss,linear,upwind,
        cellLimited,corrected,limited,vanLeer,multivariateSelection,meshWave,
        thermoType,mixture,transport,thermo,energy,equationOfState,specie,
        heRhoThermo,reactingMixture,sutherland,janaf,sensibleEnthalpy,
        perfectGas,inertSpecie,chemistryReader,foamChemistryReader,species,
        molWeight,Tlow,Thigh,Tcommon,highCpCoeffs,lowCpCoeffs,As,Ts,
        simulationType,RAS,LES,RASModel,kOmegaSST,kEpsilon,realizableKE,
        SpalartAllmaras,turbulence,printCoeffs,Sct,kMin,omegaMin,
        castellatedMesh,snap,addLayers,geometry,triSurfaceMesh,
        searchableSphere,searchableCylinder,castellatedMeshControls,
        maxLocalCells,maxGlobalCells,minRefinementCells,maxLoadUnbalance,
        nCellsBetweenLevels,features,refinementSurfaces,refinementRegions,
        resolveFeatureAngle,locationInMesh,allowFreeStandingZoneFaces,
        snapControls,nSmoothPatch,nSolveIter,nRelaxIter,nFeatureSnapIter,
        implicitFeatureSnap,explicitFeatureSnap,multiRegionFeatureSnap,
        addLayersControls,relativeSizes,layers,nSurfaceLayers,
        expansionRatio,finalLayerThickness,minThickness,nGrow,featureAngle,
        slipFeatureAngle,nSmoothSurfaceNormals,nSmoothNormals,
        nSmoothThickness,maxFaceThicknessRatio,maxThicknessToMedialRatio,
        minMedialAxisAngle,nBufferCellsNoExtrude,nLayerIter,
        meshQualityControls,maxNonOrtho,maxBoundarySkewness,
        maxInternalSkewness,maxConcave,minVol,minTetQuality,minArea,
        minTwist,minDeterminant,minFaceWeight,minVolRatio,minTriangleTwist,
        nSmoothScale,errorReduction,writeFlags,scalarLevels,mergeTolerance,
        scale,xMin,xMax,yMin,yMax,zMin,zMax,vertices,blocks,hex,
        simpleGrading,edges,boundary,faces,patch,allBoundary,
        mergePatchPairs,numberOfSubdomains,method,scotch,hierarchical,
        residuals,solverInfo,yPlus,yPlusAudit,surfaceFieldValue,
        sampledSurfaceDict,plane,planeType,pointAndNormal,pointAndNormalDict,
        point,normal,interpolate,areaAverage,sum,fieldAverage,timeStart,
        mean,prime2Mean,base,time,timeStep,runTime,writeTime,
        nonInterpolated,interpolated,levels,inside,outside,patchInfo,
        potentialFlow,transonic,fixedFluxPressure,prghTotalPressure,p,
        p_rgh,U,T,H2,CH4,h,he,K,Yi,Yi_h,k,omega,epsilon,nut,alphat,phi,
        true,false,yes,no,on,off
    },
    sensitive=true,
    morecomment=[l]{//},
    morecomment=[s]{/*}{*/},
    morestring=[b]",
}
\lstdefinestyle{foamstyle}{language=foam,basicstyle=\ttfamily\scriptsize}
\lstdefinestyle{bashstyle}{language=bash,basicstyle=\ttfamily\scriptsize}
\lstdefinestyle{cppstyle}{language=C++,basicstyle=\ttfamily\scriptsize}
\lstdefinestyle{pystyle}{language=Python,basicstyle=\ttfamily\scriptsize}

% Boxes without external packages: a coloured left bar + indented paragraph.
\newenvironment{vivabox}[1][]{%
    \par\medskip\noindent
    \begin{list}{}{%
        \setlength{\leftmargin}{1.4em}%
        \setlength{\rightmargin}{0pt}%
        \setlength{\topsep}{0pt}%
    }\item[]%
    \textcolor{orange!60!black}{\rule[-0.4em]{2.5pt}{1.0em}}%
    \hspace{0.6em}{\sffamily\bfseries\color{orange!50!black}#1}\par\smallskip
}{%
    \end{list}\medskip
}
\newenvironment{whybox}[1][]{%
    \par\medskip\noindent
    \begin{list}{}{%
        \setlength{\leftmargin}{1.4em}%
        \setlength{\rightmargin}{0pt}%
        \setlength{\topsep}{0pt}%
    }\item[]%
    \textcolor{blue!50!black}{\rule[-0.4em]{2.5pt}{1.0em}}%
    \hspace{0.6em}{\sffamily\bfseries\color{blue!50!black}#1}\par\smallskip
}{%
    \end{list}\medskip
}

\usepackage[hidelinks]{hyperref}

% Shorthand for the dimensionless groups used throughout the document.
\newcommand{\Hi}{\ensuremath{\mathrm{H}_2}}
\newcommand{\CHi}{\ensuremath{\mathrm{CH}_4}}
\newcommand{\dD}{\ensuremath{d/D}}
\newcommand{\VR}{\ensuremath{\mathrm{VR}}}
\newcommand{\HBR}{\ensuremath{\mathrm{HBR}}}
\newcommand{\CoV}{\ensuremath{\mathrm{CoV}}}
\newcommand{\Ri}{\ensuremath{\mathrm{Ri}}}
\newcommand{\Frnum}{\ensuremath{\mathrm{Fr}}}

\title{\sffamily\bfseries OpenFOAM and the H$_2$/CH$_4$ T-Junction Project --- Self-Study Notes for the MTP Viva}
\author{Sakhare Yash Balram (2022CH71496)\\
Department of Chemical Engineering, IIT Delhi}
\date{April 2026}

\begin{document}
\maketitle

\begin{abstract}
\noindent
This document is a single-stop reference for everything I need to be
able to discuss in the viva: how OpenFOAM works as a CFD framework,
what is in every file of my case directory, why each numerical and
physical choice was made, and what the project pipeline looks like
end-to-end. It is dense on purpose --- it is meant to read like
lecture notes, not a polished report. Equations and code listings are
shown in full where they would otherwise force me to wave my hands.
\end{abstract}

\tableofcontents
\clearpage

% =====================================================================
\section{What OpenFOAM actually is}
% =====================================================================

OpenFOAM (\emph{Open-source Field Operation And Manipulation}) is a
C\texttt{++} library and a command-line toolbox for solving systems of
PDEs on unstructured meshes by the \textbf{Finite Volume Method
(FVM)}. It is not a single black-box solver: it is a collection of
$\sim 100$ pre-built executables (mesh utilities, solvers,
post-processors) plus a library you can extend.

The three things to remember:

\begin{enumerate}[leftmargin=1.5em]
    \item \textbf{It is a library, not a GUI.}\;
    Every case is a directory of plain-text dictionaries. There is no
    ``project file''. You set up the case by writing dictionaries; you
    run the case by calling executables in sequence; you visualise
    results in ParaView (or by post-processing scripts).
    \item \textbf{It uses the Finite Volume Method.}\;
    The PDEs (Navier--Stokes, species transport, $k$, $\omega$, $h$,
    $\rho$\dots) are integrated over each finite cell; surface fluxes
    are reconstructed from cell-centre values via interpolation
    schemes; boundary conditions are applied face-by-face.
    \item \textbf{Solvers are application-specific.}\;
    Different physics gets different solvers --- e.g.\ \texttt{simpleFoam}
    (steady incompressible), \texttt{pimpleFoam} (transient
    incompressible), \texttt{rhoPimpleFoam} (transient compressible),
    \texttt{rhoReactingFoam} (compressible with species), and
    \texttt{rhoReactingBuoyantFoam} (compressible $+$ species $+$
    gravity, which is the one we use).
\end{enumerate}

\begin{whybox}[Why OpenFOAM and not ANSYS Fluent?]
Three reasons that hold up in viva:
\begin{enumerate}[leftmargin=1.5em,topsep=2pt]
    \item \emph{Cost / licensing.} A Fluent licence is not feasible to
    keep alive on a personal workstation for a year-long DoE; OpenFOAM
    is free and runs on Linux/WSL out of the box.
    \item \emph{Reproducibility.} The case is plain text and lives in
    git; the entire DoE can be regenerated from a seed and a stamp
    script. With Fluent the journal/.cas binary is hard to diff or
    review.
    \item \emph{Extensibility.} Custom boundary conditions
    (\texttt{codedFixedValue} 1/7-th power-law, our angle-tilted
    branch profile) drop straight into the case as inline C\texttt{++}
    snippets, no UDF compilation step.
\end{enumerate}
\end{whybox}

\subsection{The Finite Volume Method in one paragraph}

For each conservation law $\partial_t \phi + \nabla\!\cdot\!\mathbf{F}(\phi) = S(\phi)$,
the FVM integrates over a control volume $V_P$ and uses the
divergence theorem:
\begin{equation}
\frac{d}{dt}\int_{V_P}\phi \,dV \;+\; \oint_{\partial V_P}\mathbf{F}\!\cdot\!\hat{\mathbf{n}}\,dA \;=\; \int_{V_P} S \,dV.
\end{equation}
The cell-centred $\phi_P$ is the unknown; surface integrals are
discretised by summing over faces $f$ of $P$:
\begin{equation}
V_P \frac{\phi_P^{n+1}-\phi_P^{n}}{\Delta t}
\;+\; \sum_{f \in \partial V_P} \mathbf{F}_f \!\cdot\! \mathbf{S}_f
\;=\; V_P S_P.
\end{equation}
The flux $\mathbf{F}_f$ at the face needs $\phi_f$, which is interpolated
from $\phi_P$ and $\phi_N$ (neighbour) by an \emph{interpolation
scheme}: linear, upwind, vanLeer, MUSCL\dots\ The choice of scheme is
what \texttt{system/fvSchemes} controls.

\subsection{The OpenFOAM solver lifecycle}

Every CFD case in OpenFOAM goes through the same sequence:

\begin{enumerate}[leftmargin=1.5em]
    \item \textbf{Geometry.} Either generate a mesh directly with
    \texttt{blockMesh} (block-structured) or build it from STLs with
    \texttt{snappyHexMesh}.
    \item \textbf{Initial and boundary conditions.} Edit the
    dictionaries in \texttt{0/} (one per field).
    \item \textbf{Numerics.} Edit \texttt{system/fvSchemes} (spatial
    schemes), \texttt{system/fvSolution} (linear solvers,
    pressure-velocity coupling), \texttt{system/controlDict}
    (time stepping, write controls, function objects).
    \item \textbf{Decomposition.} For parallel runs,
    \texttt{decomposePar} splits the mesh into $N$ sub-meshes
    (\texttt{processor0/}, \texttt{processor1/}, \dots).
    \item \textbf{Solve.} \texttt{mpirun -np N solver -parallel} runs
    the actual transient.
    \item \textbf{Reconstruct.} \texttt{reconstructPar} merges the
    decomposed time directories back into a single \texttt{0.1/},
    \texttt{0.2/}, \dots, on the host process.
    \item \textbf{Post-process.} ParaView, function objects (sampled
    planes, surface integrals), or external scripts for derived
    quantities.
\end{enumerate}

\clearpage
% =====================================================================
\section{Anatomy of an OpenFOAM case directory}
% =====================================================================

Every OpenFOAM case --- ours included --- has exactly three top-level
directories, and that's the contract.

{\small
\begin{verbatim}
case_03/
|
+-- 0/                      <-- initial + boundary conditions per field
|   |-- U                       velocity
|   |-- p                       absolute pressure
|   |-- p_rgh                   modified pressure (= p - rho*g*x)
|   |-- T                       temperature
|   |-- H2, CH4                 species mass fractions
|   |-- k, omega                turbulence quantities
|   `-- nut, alphat             eddy-viscosity / -diffusivity
|
+-- constant/               <-- physical models (do not change with time)
|   |-- polyMesh/               the actual mesh (created by snappy/blockMesh)
|   |-- triSurface/             the STL geometry  (input to snappyHexMesh)
|   |-- thermophysicalProperties  thermo + species + transport
|   |-- turbulenceProperties      RANS model selection
|   |-- momentumTransport         alias of the above (forward-compat)
|   |-- reactions                 (empty -- no combustion)
|   |-- thermo.compressibleGas    (empty placeholder)
|   |-- combustionProperties      "off"
|   `-- g                        gravitational acceleration vector
|
`-- system/                 <-- "how to solve it"
    |-- controlDict          time-stepping, write controls, function objects
    |-- fvSchemes            ddt / div / grad / laplacian schemes
    |-- fvSolution           linear solvers + PIMPLE settings
    |-- blockMeshDict        background mesh definition
    |-- snappyHexMeshDict    snapped-Cartesian-cut-cell controls
    |-- decomposeParDict     parallel domain decomposition
    |-- createPatchDict      patch renaming / merging
    |-- surfaceFeatureExtractDict   feature edge extraction for snappy
    `-- postProcess.dict     post-processing recipes (optional)
\end{verbatim}
}

\paragraph{Why three directories?}
The split mirrors the three independent decisions in any CFD problem:
\emph{what is the state at $t=0$ and on the boundaries} (\texttt{0/}),
\emph{what physical model do we use} (\texttt{constant/}), and
\emph{what numerics do we use to solve it} (\texttt{system/}). They
are physically separate so e.g.\ a mesh study can swap
\texttt{constant/polyMesh/} without touching the BCs, and a
schemes study can swap \texttt{system/fvSchemes} without touching the
mesh.

\clearpage
% =====================================================================
\section{The \texttt{system/} directory --- numerics and controls}
% =====================================================================

\subsection{controlDict --- time, output, diagnostics}
\label{sec:controlDict}

\texttt{controlDict} controls \emph{when} the solver runs, \emph{how
often} it writes, and \emph{what runtime diagnostics} it computes.
Annotated extract from our case~03:

\begin{lstlisting}[style=foamstyle,caption={\texttt{system/controlDict} (annotated)}]
application         rhoReactingBuoyantFoam;   // <-- which solver to use

startFrom           startTime;                // start from t = startTime
startTime           0;                        //   (latestTime to resume)
stopAt              endTime;                  // stop at t = endTime
endTime             1.2;                      // 1.2 s = ~2 flow-throughs
deltaT              1e-5;                     // initial dt; will adapt

writeControl        adjustable;               // write at runTime intervals
writeInterval       0.1;                      // write every 0.1 s of sim time
purgeWrite          3;                        // keep last 3 time-dirs
writePrecision      10;
runTimeModifiable   true;                     // re-read controlDict on the fly

adjustTimeStep      yes;                      // CFL-based dt adaptation
maxCo               4.5;                      // CFL cap
maxDeltaT           0.01;                     // hard cap on dt

functions
{
    residuals { type solverInfo; ... }       // per-step residual log
    yPlusAudit { type yPlus; ... }           // y+ stats every 0.1s
    mainInletFlux  { type surfaceFieldValue; fields (phi); ... }
    branchInletFlux{ type surfaceFieldValue; fields (phi); ... }
    outletFlux     { type surfaceFieldValue; fields (phi); ... }
    H2_outletAvg   { type surfaceFieldValue; fields (H2); operation areaAverage; }
    p_rgh_inlet    { ...areaAverage of p_rgh on main_inlet... }
    p_rgh_outlet   { ...areaAverage of p_rgh on outlet... }
    plane_z3_H2    { type surfaceFieldValue; regionType sampledSurface;
                     sampledSurfaceDict { type plane; ... } }
    // similar for z=4,5,6
    fieldAverage   { type fieldAverage; timeStart 0.6;
                     fields ( U  H2  CH4  p_rgh  k ); }
}
\end{lstlisting}

\subsubsection{Things to know about controlDict}

\begin{itemize}[leftmargin=1.5em]
    \item \textbf{\texttt{adjustTimeStep yes; maxCo 4.5}} : the solver
    automatically picks $\Delta t$ at every step so that the maximum
    Courant number $\mathrm{Co}=|u|\,\Delta t/\Delta x$ across the mesh
    stays below 4.5. PIMPLE with two correctors can stably handle
    $\mathrm{Co}\sim 5$. We later raised this to 8.0 for the 30\textdegree{}
    DoE to halve wall time.
    \item \textbf{\texttt{writeInterval 0.1}} : the solver writes a
    full snapshot every 0.1\,s of simulated time. Combined with
    \texttt{purgeWrite 3}, only the last 3 snapshots are kept on disk
    (otherwise a 1.2\,s run would write 12 directories of $\sim 30$\,MB
    each).
    \item \textbf{Function objects} run \emph{inside the solver loop},
    not as a post-processing pass. They are the right way to compute
    $\Delta p$, $y^+$, and area-averaged $H_2$ at every step without
    storing every snapshot.
    \item \textbf{\texttt{fieldAverage}} starts time-averaging at
    $t_\mathrm{start}=0.6$\,s, after the start-up transient has cleared
    (one main-pipe flow-through). \texttt{prime2Mean on} for $H_2$ also
    accumulates the second moment so we can compute $\sigma_{H_2}^2$.
\end{itemize}

\subsection{fvSchemes --- spatial and temporal schemes}

\texttt{fvSchemes} tells the solver which numerical scheme to use for
each derivative term in the equations.

\begin{lstlisting}[style=foamstyle,caption={\texttt{system/fvSchemes} (full)}]
ddtSchemes
{
    default         Euler;                          // 1st-order implicit
}

gradSchemes
{
    default         Gauss linear;
    grad(U)         cellLimited Gauss linear 1;     // limited to bound velocity
    grad(p_rgh)     Gauss linear;
    grad(H2)        cellLimited Gauss linear 1;
}

divSchemes
{
    default         none;                           // force every term to be set

    div(phi,U)      Gauss upwind;                   // 1st-order upwind, stable
    div(phi,H2)     Gauss upwind;
    div(phi,CH4)    Gauss upwind;
    div(phi,k)      Gauss upwind;
    div(phi,omega)  Gauss upwind;
    div(phi,h)      Gauss upwind;
    div(phi,K)      Gauss upwind;
    div(phi,Yi_h)   Gauss multivariateSelection { H2 upwind; CH4 upwind; h upwind; };
    div(phiv,p)     Gauss upwind;
    div(((rho*nuEff)*dev2(T(grad(U))))) Gauss linear;   // viscous stress
}

laplacianSchemes
{
    default         Gauss linear corrected;          // 2nd-order, snGrad-corrected
}

interpolationSchemes
{
    default         linear;                          // central interpolation
}

snGradSchemes
{
    default         corrected;                       // surface-normal grad with non-orthog correction
}

wallDist
{
    method          meshWave;                        // wall distance for SST
}

fluxRequired
{
    default         no;
    p_rgh;                                           // p_rgh flux is needed for momentum
}
\end{lstlisting}

\begin{whybox}[Why first-order upwind for div?]
Upwind is \emph{numerically diffusive} (it artificially smears sharp
gradients), but it is \emph{monotonic} --- it does not generate
wiggles. For $H_2$ mass fraction this matters a lot because second-
or higher-order schemes can overshoot beyond $[0,1]$ in cells with
sharp gradients (the jet edge), which the JANAF thermophysics will
then reject as unphysical. The cost is $\sim 5$--$10\%$ excess
diffusion in CoV (Ferziger \& Peri\'c 2002, sec.\ 8.8), which is
acceptable for a parametric DoE where we look at \emph{trends} not
absolute numbers. Compare with literature: Su et al.\ also use
first/second-order upwind for $k$, $\varepsilon$ and energy.
\end{whybox}

\subsubsection{What ``Gauss linear corrected'' means}

\begin{itemize}[leftmargin=1.5em,topsep=2pt]
    \item \textbf{Gauss}: surface-integral form (the default for FVM).
    \item \textbf{linear}: linear interpolation from cell centres to
    faces (i.e.\ $\phi_f = w\phi_P + (1-w)\phi_N$).
    \item \textbf{corrected}: the surface-normal gradient at $f$ is
    corrected for mesh non-orthogonality (i.e.\ when $P$, $N$, and the
    face centre are not collinear). This is what prevents
    second-order schemes from breaking down on the kind of polyhedral
    cells that snappyHexMesh produces.
\end{itemize}

\subsection{fvSolution --- linear solvers and PIMPLE}

\texttt{fvSolution} is a two-part dictionary:

\begin{enumerate}[leftmargin=1.5em,topsep=2pt]
    \item \texttt{solvers \{...\}} : how to solve the linear system
    associated with each field.
    \item \texttt{PIMPLE \{...\}} : how the
    pressure--velocity coupling loop is structured.
\end{enumerate}

\begin{lstlisting}[style=foamstyle,caption={\texttt{system/fvSolution} (annotated extract)}]
solvers
{
    "rho.*"     { solver diagonal; }                     // density: trivial diagonal solve

    p           { solver GAMG; smoother GaussSeidel;     // GAMG = geometric algebraic multigrid
                  tolerance 1e-7; relTol 0.01; }

    p_rgh       { solver PCG; preconditioner DIC;        // PCG = preconditioned conjugate gradient
                  tolerance 1e-7; relTol 0.01; maxIter 200; }
    p_rghFinal  { $p_rgh; relTol 0; }                    // last corrector: solve to convergence

    "U"         { solver smoothSolver; smoother symGaussSeidel;
                  tolerance 1e-7; relTol 0.1; nSweeps 2; }
    UFinal      { $U; relTol 0; }

    "(H2|CH4|Yi|k|omega|epsilon)"
                { solver PBiCGStab; preconditioner DILU;
                  tolerance 1e-7; relTol 0.1; }
    "(H2|CH4|Yi|k|omega|epsilon)Final" { $H2; relTol 0; }
}

PIMPLE
{
    momentumPredictor        yes;     // step (a): predict U from old p
    nOuterCorrectors         1;       // # of full PISO outer loops per dt
    nCorrectors              2;       // # of pressure corrections per outer
    nNonOrthogonalCorrectors 1;       // # of mesh-non-ortho corrections
    transonic                no;
}

relaxationFactors
{
    equations { "U.*" 0.7; "k.*" 0.7; "omega.*" 0.7;
                "H2.*" 0.9; "CH4.*" 0.9; }
    fields    { "p_rgh.*" 0.3; }
}
\end{lstlisting}

\begin{whybox}[Why PCG/DIC for $p_\mathrm{rgh}$ instead of GAMG?]
GAMG (geometric--algebraic multigrid) is normally faster, but on
compressible flows with steep density contrasts at the junction, it
sometimes triggers a floating-point exception in the agglomeration
step (the coarse-level coefficients become ill-conditioned). PCG
with DIC preconditioner is slower per iteration but rock-solid;
for our $\sim 460$\,k cell mesh it converges in
$\lesssim 50$ iterations, which is fine.
\end{whybox}

\subsubsection{What PIMPLE actually does, step by step}

PIMPLE = \emph{merged PISO\,+\,SIMPLE}. For each timestep:

\begin{enumerate}[leftmargin=1.5em,topsep=2pt]
    \item \textbf{Momentum predictor} (if \texttt{momentumPredictor
    yes}): solve momentum with old $p$, get a tentative $\mathbf{u}^*$.
    \item \textbf{Pressure correction} (PISO loop, repeated
    \texttt{nCorrectors} times):
    \begin{itemize}[leftmargin=1em,topsep=1pt]
        \item Build the pressure equation from continuity using
        $\mathbf{u}^*$.
        \item Solve $p_{\mathrm{rgh}}$.
        \item Correct $\mathbf{u}$ using the new pressure gradient.
    \end{itemize}
    \item \textbf{Repeat} (SIMPLE outer loop, \texttt{nOuterCorrectors}
    times) if outer-loop convergence is needed within the timestep.
    \item \textbf{Solve scalars}: $k$, $\omega$, $h$, $H_2$.
    \item \textbf{Update $\rho$, $\mu$, $\nu_t$} from thermo and
    turbulence model.
\end{enumerate}

For our \texttt{nOuterCorrectors=1, nCorrectors=2}, that is
``one outer iteration with two PISO corrections per timestep'' --- the
PISO mode of PIMPLE, suitable for transient flows where the time step
is small enough (\(\mathrm{Co}\le 5\)).

\begin{vivabox}[Likely viva question]
\textbf{``Why one outer corrector and not three?''}\\
Because we are doing a \emph{transient} simulation with a
small time step (Co $\le 4.5$). In transient mode, the field changes
between time steps are small, so a single outer iteration plus two
PISO inner corrections is enough to drop the residuals by 5--6 orders
of magnitude per step. More outer iterations would just multiply wall
time without improving accuracy --- they only help when you push
$\mathrm{Co}\gg 1$ to chase a steady state.
\end{vivabox}

\subsection{blockMeshDict --- the background mesh}

\texttt{blockMesh} reads \texttt{system/blockMeshDict} and emits a
structured hex mesh of the bounding box. This is the \emph{background}
mesh that snappyHexMesh later carves the geometry out of.

\begin{lstlisting}[style=foamstyle,caption={\texttt{system/blockMeshDict} (key parts)}]
scale   1;

xMin 0.00; xMax 0.30;     // half-domain (x = 0 is symmetry plane)
yMin -0.30; yMax 1.50;
zMin -0.05; zMax 6.95;    // R2 shortened upstream length (was up to z=-2.3)

vertices ( ($xMin $yMin $zMin) ... ($xMin $yMax $zMax) );  // 8 corners

blocks
(
    hex (0 1 2 3 4 5 6 7) (10 60 232) simpleGrading (1 1 1)
);
//                       ^^^^^^^^^^^^
//   uniform 30 mm cells in all 3 directions = 139,200 background cells

boundary
(
    sym
    {
        type symmetryPlane;             // x = 0 face is the mirror plane
        faces ( (0 4 7 3) );
    }
    allBoundary
    {
        type patch;                     // catch-all for snappy to retag
        faces ( ... five other faces ... );
    }
);
\end{lstlisting}

\paragraph{Two things to flag.}
First, \texttt{xMin = 0} means we run only the half-domain
$x\ge 0$ and recover the other half by mirror symmetry, halving the
cell count. Second, \texttt{allBoundary} is a deliberate catch-all
for the other five faces of the bounding box --- snappyHexMesh will
later split it into the actual \texttt{wall}, \texttt{main\_inlet},
\texttt{outlet} and \texttt{branch\_inlet} patches based on the STL
files.

\subsection{snappyHexMeshDict --- carving the geometry}

snappyHexMesh has three phases controlled by three flags at the top:

\begin{lstlisting}[style=foamstyle]
castellatedMesh true;     // (1) hex-cut cells where they intersect STLs
snap            true;     // (2) snap cell points to STL surface
addLayers       true;     // (3) extrude prismatic layers along walls
\end{lstlisting}

\subsubsection{Phase 1: castellation (refinement)}

\begin{lstlisting}[style=foamstyle,caption={\texttt{snappyHexMeshDict} -- castellation}]
geometry
{
    wall.stl          { type triSurfaceMesh; name wall; }
    main_inlet.stl    { type triSurfaceMesh; name main_inlet; }
    outlet.stl        { type triSurfaceMesh; name outlet; }
    branch_inlet.stl  { type triSurfaceMesh; name branch_inlet; }

    // Sphere of refinement around the junction
    junctionRefine
    {
        type        searchableSphere;
        centre      (0 0 2.3);                  // junction centre
        radius      0.69;                       // 1.5*D_main
    }
    // Cylindrical refinement extending downstream of junction
    jetRefine
    {
        type        searchableCylinder;
        point1      (0 0 2.6);                  // just past junction
        point2      (0 0 6.2);                  // ~8 D_main downstream
        radius      0.23;                       // R_main
    }
}

castellatedMeshControls
{
    maxLocalCells       200000;
    maxGlobalCells      1500000;
    minRefinementCells  10;
    nCellsBetweenLevels 3;

    features                                    // edge refinement from .eMesh files
    (
        { file "wall.eMesh";         level 1; }
        { file "main_inlet.eMesh";   level 0; }
        { file "outlet.eMesh";       level 0; }
        { file "branch_inlet.eMesh"; level 1; }
    );

    refinementSurfaces
    {
        wall          { level (1 1); patchInfo { type wall; } }   // 1 level of surface refinement
        main_inlet    { level (0 0); patchInfo { type patch; } }
        outlet        { level (0 0); patchInfo { type patch; } }
        branch_inlet  { level (1 1); patchInfo { type patch; } }
    }

    resolveFeatureAngle 30;

    refinementRegions
    {
        junctionRefine { mode inside; levels ((1e15 2)); }   // level 2 inside sphere
        jetRefine      { mode inside; levels ((1e15 1)); }   // level 1 inside cylinder
    }

    locationInMesh (0.077 0.011 1.111);    // any point physically inside the fluid
    allowFreeStandingZoneFaces true;
}
\end{lstlisting}

The way to read \texttt{level (1 1)} : a refinement level $L$ means
the cell side length is $\frac{1}{2^L}$ of the background cell. Our
background is 30\,mm, so \texttt{level 1} $\rightarrow$ 15\,mm, and the
inside of the junction sphere at \texttt{level 2} $\rightarrow$ 7.5\,mm.

\subsubsection{Phase 2: snap and Phase 3: layer addition}

\begin{lstlisting}[style=foamstyle]
snapControls
{
    nSmoothPatch            3;
    tolerance               2.0;
    nSolveIter              50;
    nRelaxIter              5;
    nFeatureSnapIter        10;
    explicitFeatureSnap     true;
}

addLayersControls
{
    relativeSizes        true;          // layer thickness as fraction of adjacent cell
    layers
    {
        wall { nSurfaceLayers 3; }      // 3 prism layers on every wall face
    }
    expansionRatio       1.2;           // each layer 1.2x thicker than the one below
    finalLayerThickness  0.5;           // outermost layer = 0.5 of adjacent cell
    minThickness         0.15;          // give up if a layer would be < 15%
    nLayerIter           40;
}
\end{lstlisting}

The \emph{snap} phase moves cell points onto the STL surface so the
mesh actually conforms to the geometry. The \emph{addLayers} phase
then extrudes 3 prismatic layers from the wall inward to land
$y^+\sim 30$--$80$ at the first cell, which is the wall-function
regime our $k$--$\omega$ SST setup needs.

\begin{vivabox}[Likely viva question]
\textbf{``Why $y^+\sim 30$--$80$ instead of resolving to $y^+\sim 1$?''}\\
With $y^+\sim 1$ we would need $\sim 5\times$ more cells in the
boundary layer and the time step would have to drop by a similar
factor, since the smallest cell sets the CFL limit. For a parametric
DoE this is unaffordable. Wall functions
(\texttt{kqRWallFunction}, \texttt{omegaWallFunction}) bridge the
viscous sublayer analytically, which is industry standard for the
Reynolds numbers we have ($Re\sim 5\times 10^5$) and the
quantities we report (CoV, $\Delta p$). For LES or
heat-transfer studies this trade-off would not be acceptable.
\end{vivabox}

\subsection{decomposeParDict --- parallel decomposition}

\begin{lstlisting}[style=foamstyle,caption={\texttt{system/decomposeParDict} (full)}]
numberOfSubdomains  16;
method              scotch;
\end{lstlisting}

\texttt{scotch} is a graph-partitioning library that minimises the
inter-process communication surface. The mesh is split into 16
sub-meshes (16 MPI ranks on the 8-core/16-thread workstation), each
running a copy of the solver on its own \texttt{processor*/} folder.
Each rank only needs values from neighbour ranks at sub-domain
boundary faces (the ``halo''), which is what scotch tries to keep
small.

\clearpage
% =====================================================================
\section{The \texttt{constant/} directory --- physical models}
% =====================================================================

\subsection{turbulenceProperties (and momentumTransport)}

\begin{lstlisting}[style=foamstyle,caption={\texttt{constant/turbulenceProperties} (full)}]
simulationType  RAS;        // RAS = Reynolds-Averaged Simulation (RANS)

RAS
{
    RASModel        kOmegaSST;
    turbulence      on;
    printCoeffs     on;     // dump model constants to log on first call
}

Sct     0.5;                // Turbulent Schmidt number for species transport
kMin    1e-10;              // floor on k to avoid divide-by-zero
omegaMin 1e-10;
\end{lstlisting}

\paragraph{Schmidt number $Sc_t = 0.5$.}
This sets the turbulent diffusivity for species
($D_\mathrm{turb} = \nu_t / Sc_t$). The textbook range is 0.7--1.0
for inert gases; 0.5 is the calibrated value for high-density-ratio
mixing (Su 2022, He et al.\ 2019).

\subsubsection{The k-$\omega$ SST model in 30 seconds}

\begin{itemize}[leftmargin=1.5em,topsep=2pt]
    \item Solves two transport equations: one for $k$ (turbulent
    kinetic energy) and one for $\omega$ (specific dissipation rate
    $\omega = \varepsilon/k$).
    \item Eddy viscosity: $\mu_t = \rho \, a_1\, k / \max(a_1\,\omega,\,F_2 S)$
    where $a_1=0.31$, $F_2$ is a blending function active near walls,
    $S$ is strain-rate magnitude. The $\max(\cdot)$ is the
    \emph{shear-stress transport} (SST) limiter.
    \item Blending function $F_1$ switches between standard
    $k$--$\omega$ near walls (where it's good) and standard
    $k$--$\varepsilon$ in the free shear (where it's better).
    \item Wall functions in our case
    (\texttt{kqRWallFunction}/\texttt{omegaWallFunction}) bridge the
    viscous sublayer analytically when $y^+ > 11$.
\end{itemize}

\begin{whybox}[Why SST and not standard k-$\varepsilon$ (which Su 2022 used)?]
\begin{enumerate}[leftmargin=1.5em,topsep=2pt]
    \item \textbf{Adverse pressure gradients.} $k$--$\varepsilon$ over-predicts
    eddy viscosity in adverse $\nabla p$ regions (the lee side of the
    branch jet). SST limits this through $F_2$.
    \item \textbf{Separation prediction.} $k$--$\varepsilon$ under-predicts
    separation; SST handles separation length within $\sim 5$\%.
    \item \textbf{Wall behaviour.} $k$--$\varepsilon$ needs damping
    functions or a two-layer model near walls; SST does not.
\end{enumerate}
The penalty is solving two equations whose coupling is more sensitive
to discretisation. We pay that with cell-limited gradients in
\texttt{fvSchemes} and PBiCGStab+DILU on $k,\omega$.
\end{whybox}

\subsection{thermophysicalProperties}

This is where multispecies thermo lives.

\begin{lstlisting}[style=foamstyle,caption={\texttt{constant/thermophysicalProperties} (key parts)}]
thermoType
{
    type            heRhoThermo;        // density-based, sensible enthalpy
    mixture         reactingMixture;    // multi-species mixture
    transport       sutherland;         // Sutherland's law for mu(T)
    thermo          janaf;              // JANAF polynomials for cp(T)
    energy          sensibleEnthalpy;   // h_sensible = integral(cp dT)
    equationOfState perfectGas;         // p = rho R T
    specie          specie;
}

inertSpecie      CH4;                  // not solved for; closure species
species ( H2 CH4 );

H2
{
    specie { molWeight 2.016; }        // kg/kmol
    thermodynamics
    {
        Tlow 200; Thigh 6000; Tcommon 1000;
        highCpCoeffs ( 2.93287 8.27e-04 -1.46e-07 1.54e-11 -6.89e-16 -813.07 -1.024 );
        lowCpCoeffs  ( 2.34433 7.98e-03 -1.95e-05 2.02e-08 -7.38e-12 -917.94 0.683 );
    }
    transport { As 6.69e-07; Ts 72.0; }   // Sutherland mu = As * sqrt(T) / (1 + Ts/T)
}

CH4
{
    specie { molWeight 16.043; }
    thermodynamics { ... 7-coeff JANAF polynomial ... }
    transport      { As 1.17e-06; Ts 164.0; }
}
\end{lstlisting}

\paragraph{What \texttt{heRhoThermo} buys us.}
\texttt{heRhoThermo} is a density-based thermodynamics class:
$\rho$ is computed from $p$, $T$ and the local mixture composition,
and the energy variable is sensible enthalpy
$h = \sum_k Y_k \int_{T_\mathrm{ref}}^{T} c_{p,k}(T') \,dT'$. JANAF
polynomials provide $c_p(T)$ as a 7-coefficient piecewise polynomial
split at $T_\mathrm{common}=1000$\,K. This is what lets us treat
$H_2/CH_4$ as a real binary mixture rather than two passive scalars.

\subsection{The other small files in constant/}

\begin{lstlisting}[style=foamstyle,caption={\texttt{constant/g} (full)}]
dimensions  [0 1 -2 0 0 0 0];
value       (0 -9.81 0);                     // gravity points -y (down)
\end{lstlisting}

\noindent
This is what activates the buoyancy term in
\texttt{rhoReactingBuoyantFoam}. \texttt{combustionProperties} is set to
``off'' (we don't want chemistry), \texttt{reactions} is empty, and
\texttt{thermo.compressibleGas} is a placeholder demanded by the
\texttt{foamChemistryReader} format.

\clearpage
% =====================================================================
\section{The \texttt{0/} directory --- initial and boundary conditions}
% =====================================================================

There is one dictionary per field variable. The structure is always
the same: \texttt{dimensions}, \texttt{internalField}, then a
\texttt{boundaryField} block listing one entry per patch.

\subsection{0/U --- velocity (the most interesting one)}

This is where the \emph{custom} 1/7-th power-law profile lives.
Below is the full \texttt{codedFixedValue} for the main inlet and
the branch inlet at $\alpha=90^\circ$ from case~03:

\begin{lstlisting}[style=foamstyle,caption={\texttt{0/U} -- 90\textdegree{} version (case~03)}]
internalField   uniform (0 0 5);

boundaryField
{
    wall { type noSlip; }
    sym  { type symmetryPlane; }

    main_inlet
    {
        type            codedFixedValue;
        value           uniform (0 0 10);
        name            mainInletPowerLaw;
        code
        #{
            const vectorField& Cf = this->patch().Cf();   // face centres
            vectorField& field = *this;                   // the BC values
            const scalar d    = 0.46;                     // main pipe D
            const scalar uMax = 10;                       // main bulk velocity
            const scalar n    = 7.0;                      // 1/7-th law exponent
            forAll(Cf, i)
            {
                const scalar r = Foam::sqrt(sqr(Cf[i].x()) + sqr(Cf[i].y()));
                const scalar arg = Foam::max(1.0 - sqr(2.0*r/d), scalar(0));
                const scalar u   = uMax * Foam::pow(arg, 1.0/n);
                field[i] = vector(0, 0, u);
            }
        #};
    }

    outlet { type pressureInletOutletVelocity; value uniform (0 0 5); }

    branch_inlet
    {
        type            codedFixedValue;
        value           uniform (0 -13.3032 0);
        name            branchInletPowerLaw;
        code
        #{
            const vectorField& Cf = this->patch().Cf();
            vectorField& field = *this;
            const scalar d    = 0.11585;
            const scalar uMax = 13.3032;
            const scalar zc   = 2.3;                      // junction z
            const scalar n    = 7.0;
            forAll(Cf, i)
            {
                const scalar r = Foam::sqrt(sqr(Cf[i].x()) + sqr(Cf[i].z() - zc));
                const scalar arg = Foam::max(1.0 - sqr(2.0*r/d), scalar(0));
                const scalar u   = uMax * Foam::pow(arg, 1.0/n);
                field[i] = vector(0, -u, 0);              // injecting in -y
            }
        #};
    }
}
\end{lstlisting}

\subsubsection{Reading the codedFixedValue inline-C\texttt{++}}

\begin{itemize}[leftmargin=1.5em,topsep=2pt]
    \item \texttt{this->patch().Cf()} returns the centre coordinate of
    every face on this patch.
    \item \texttt{vectorField\& field = *this} : the BC array we are
    \emph{writing} (per-face $\mathbf{u}$).
    \item For each face $i$, we compute the radial distance $r$ from
    the pipe axis, plug it into the 1/7-th power law,
    \(u(r) = u_\mathrm{max}(1 - (2r/d)^2)^{1/7}\), and write
    \texttt{field[i]}.
    \item \texttt{Foam::max(\dots,\,0)} guards against negative
    arguments at the wall (numerical safety) before the
    \texttt{Foam::pow}, which would otherwise return NaN.
\end{itemize}

\subsubsection{The 30\textdegree{} extension --- tilted branch}

For an arbitrary angle, the branch axis is no longer aligned with $-y$,
so we need to compute the radial distance \emph{in the branch's local
frame} and project the velocity vector along $-\hat{\mathbf{s}}$ where
$\hat{\mathbf{s}}$ is the branch axis. This is the version stamped at
30\textdegree{} (\(s_\alpha = \sin\alpha = 0.5,\,c_\alpha = \cos\alpha
\approx 0.866\)):

\begin{lstlisting}[style=foamstyle,caption={\texttt{0/U} -- 30\textdegree{} branch (the angle-agnostic version)}]
branch_inlet
{
    type            codedFixedValue;
    value           uniform (0 -6.2066 10.7501);   // = u_branch * (-sin a, cos a)
    name            branchInletPowerLaw;
    code
    #{
        const vectorField& Cf = this->patch().Cf();
        vectorField& field = *this;
        const scalar d    = 0.13631;
        const scalar uMax = 12.4132;
        const scalar zc   = 2.3;                   // junction z
        const scalar yJ   = 0.23;                  // junction y (= R_main)
        const scalar sa   = 0.5;                   // sin(alpha)
        const scalar ca   = 0.866025;              // cos(alpha)
        const scalar n    = 7.0;
        forAll(Cf, i)
        {
            // (1) vector from junction centre to face centre
            const scalar wx = Cf[i].x();
            const scalar wy = Cf[i].y() - yJ;
            const scalar wz = Cf[i].z() - zc;

            // (2) project onto the branch axis  S_HAT = (0, sa, -ca)
            const scalar sProj = wy*sa + wz*(-ca);

            // (3) perpendicular component (radial in the local frame)
            const scalar ppx = wx;
            const scalar ppy = wy - sProj*sa;
            const scalar ppz = wz - sProj*(-ca);
            const scalar r   = Foam::sqrt(ppx*ppx + ppy*ppy + ppz*ppz);

            // (4) 1/7-th law
            const scalar arg = Foam::max(1.0 - sqr(2.0*r/d), scalar(0));
            const scalar u   = uMax * Foam::pow(arg, 1.0/n);

            // (5) flow direction is -S_HAT = (0, -sa, ca)
            field[i] = vector(0.0, -u*sa, u*ca);
        }
    #};
}
\end{lstlisting}

\paragraph{The point of the angle-agnostic code.}
With this construction, \texttt{stamp\_cases.py} only needs to
substitute the four scalars \texttt{(d, uMax, sa, ca)} into the
template, and the same dictionary works for $\alpha = 30^\circ$,
$90^\circ$, $150^\circ$, or anything else. At
$\alpha=90^\circ$ we have $s_\alpha=1, c_\alpha=0$, and the code
reduces algebraically to the original ``\texttt{vector(0, -u, 0)}''
formulation.

\subsection{0/p\_rgh --- modified pressure}

\texttt{p\_rgh} is the modified pressure
$p_\mathrm{rgh} = p - \rho \mathbf{g}\!\cdot\!\mathbf{x}$ that lets us
absorb the hydrostatic part out of the pressure equation. This is
\emph{the} field the buoyancy solver actually solves; \texttt{p} is
reconstructed from $p_\mathrm{rgh}$ and $\rho \mathbf{g}\!\cdot\!\mathbf{x}$
at every step.

\begin{lstlisting}[style=foamstyle,caption={\texttt{0/p\_rgh} (full)}]
dimensions  [1 -1 -2 0 0 0 0];           // [Pa]
internalField   uniform 6900000;          // 6.9 MPa  (operating pressure of NG pipeline)

boundaryField
{
    main_inlet   { type fixedFluxPressure; value uniform 6900000; }
    branch_inlet { type fixedFluxPressure; value uniform 6900000; }
    outlet       { type prghPressure; p uniform 6900000; value uniform 6900000; }
    wall         { type fixedFluxPressure; value uniform 6900000; }
    sym          { type symmetryPlane; }
}
\end{lstlisting}

\subsubsection{What \texttt{fixedFluxPressure} and \texttt{prghPressure} mean}

\begin{itemize}[leftmargin=1.5em,topsep=2pt]
    \item \texttt{fixedFluxPressure} : the pressure gradient at the
    boundary is set so that the velocity-flux at the boundary is
    consistent with the actual mass flow there. It is the
    momentum-equation--consistent counterpart of the velocity BC ---
    not a fixed value of pressure.
    \item \texttt{prghPressure} (on outlet) : sets the static pressure
    \(p\) at the outlet, then computes the corresponding
    \(p_\mathrm{rgh}\) using the local density. So the user thinks
    in terms of \(p\) (intuitive), but the solver works with
    \(p_\mathrm{rgh}\) (clean).
\end{itemize}

\subsection{0/H2 and 0/CH4 --- species mass fractions}

\begin{lstlisting}[style=foamstyle,caption={\texttt{0/H2} (full)}]
dimensions      [0 0 0 0 0 0 0];               // dimensionless
internalField   uniform 0;                      // start with zero H2 everywhere

boundaryField
{
    wall         { type zeroGradient; }
    sym          { type symmetryPlane; }
    main_inlet   { type fixedValue; value uniform 0; }   // pure CH4 enters from main
    outlet       { type inletOutlet; inletValue uniform 0; value uniform 0; }
    branch_inlet { type fixedValue; value uniform 1; }   // pure H2 enters from branch
}
\end{lstlisting}

The \texttt{inletOutlet} BC handles backflow at the outlet
gracefully: if a face has $\mathbf{u}\!\cdot\!\hat{\mathbf{n}}>0$
(outflow), it acts as zeroGradient; if it has back-flow, it injects
\texttt{inletValue}. For our cases there is essentially no back-flow,
but having this BC means a transient blip never crashes the run.

\subsection{0/k and 0/omega --- turbulence inlet BCs}

For incoming patches, the standard practice is to specify $k$ from a
turbulence intensity $I$ and $\omega$ from a mixing length $\ell$:
\begin{align}
k &= \tfrac{3}{2}(I\,U)^2, & \omega &= \frac{\sqrt{k}}{C_\mu^{1/4}\,\ell}, & C_\mu &= 0.09.
\end{align}
With $I = 5\%$ (typical for $Re>10^5$ pipe flow), $U=10$\,m/s,
$D=0.46$\,m, $\ell = 0.07\,D = 0.0322$\,m, this gives
$k_\text{main}=0.375$\,m$^2$/s$^2$ and
$\omega_\text{main}=34.72$\,1/s --- exactly the numbers in our case.

\begin{lstlisting}[style=foamstyle,caption={\texttt{0/k} (full)}]
internalField uniform 0.375;

boundaryField
{
    main_inlet   { type turbulentIntensityKineticEnergyInlet; intensity 0.05; value uniform 0.375; }
    branch_inlet { type turbulentIntensityKineticEnergyInlet; intensity 0.05; value uniform 0.663661; }
    outlet       { type inletOutlet; inletValue uniform 0.375; value uniform 0.375; }
    wall         { type kqRWallFunction; value uniform 0.375; }
    sym          { type symmetryPlane; }
}
\end{lstlisting}

\begin{lstlisting}[style=foamstyle,caption={\texttt{0/omega} (full)}]
internalField uniform 35;

boundaryField
{
    main_inlet   { type turbulentMixingLengthFrequencyInlet; mixingLength 0.0322;  value uniform 34.72; }
    branch_inlet { type turbulentMixingLengthFrequencyInlet; mixingLength 0.008109; value uniform 183.41; }
    outlet       { type inletOutlet; inletValue uniform 34.72; value uniform 34.72; }
    wall         { type omegaWallFunction; value uniform 34.72; }
    sym          { type symmetryPlane; }
}
\end{lstlisting}

\subsection{0/nut, 0/alphat, 0/T, 0/p}

\begin{itemize}[leftmargin=1.5em,topsep=2pt]
    \item \texttt{nut} : turbulent kinematic viscosity. The wall
    function uses \texttt{nutkWallFunction} to set $\nu_t$ at the
    wall from $k$ and $y^+$. Internally OpenFOAM does this
    automatically; we just supply a non-zero starting field.
    \item \texttt{alphat} : turbulent thermal diffusivity. Same
    story --- a wall function field.
    \item \texttt{T} : temperature. We run isothermal at 300\,K, so
    \texttt{T} is just a constant uniform 300 with all-walls
    \texttt{zeroGradient}; the energy equation is effectively
    deactivated (relTol 1, maxIter 0 in fvSolution).
    \item \texttt{p} : absolute pressure. Reconstructed from
    \texttt{p\_rgh} every step; \texttt{0/p} is just an initial guess
    set to the same operating pressure as \texttt{0/p\_rgh}.
\end{itemize}

\clearpage
% =====================================================================
\section{The Allrun pipeline --- what actually runs}
% =====================================================================

The \texttt{Allrun} script ties everything together. Here is the
real script for case~04 (annotated):

\begin{lstlisting}[style=bashstyle,caption={\texttt{case\_04/Allrun} -- the per-case pipeline}]
#!/usr/bin/env bash
set -eo pipefail
cd "$(dirname "$0")"

# 1. Source OpenFOAM environment
OPENFOAM_BASHRC="${OPENFOAM_BASHRC:-/usr/lib/openfoam/openfoam2406/etc/bashrc}"
set +eu
source "$OPENFOAM_BASHRC" 2>/dev/null
set -e

# 2. Per-case env (UMAIN, UBRANCH, D1, D2, ZJCT, ALPHA_DEG)
source ./case.env

# 3. If we have a same-d/D donor case, copy its mesh to skip ~30 min of meshing
if [[ "${REUSE_MESH_FROM:-}" != "" && -d "${REUSE_MESH_FROM}/constant/polyMesh" ]]; then
    cp -r "${REUSE_MESH_FROM}/constant/polyMesh"   constant/
    cp -r "${REUSE_MESH_FROM}/constant/triSurface" constant/ 2>/dev/null || true
    for p in "${REUSE_MESH_FROM}"/processor*; do
        dst="./$(basename "$p")"
        mkdir -p "$dst/constant"
        cp -r "$p/constant/polyMesh" "$dst/constant/"
    done
else
    # 4. Otherwise build mesh from scratch
    python3 generateSTL.py            | tee log.generateSTL
    python3 scripts/clip_stls.py constant/triSurface | tee log.clipStls
    surfaceFeatureExtract             | tee log.surfaceFeatureExtract
    blockMesh                          | tee log.blockMesh
    snappyHexMesh -overwrite           | tee log.snappyHexMesh
    if grep -q "allBoundary" constant/polyMesh/boundary; then
        createPatch -overwrite         | tee log.createPatch
    fi
    rm -f 0/cellLevel 0/pointLevel 0/cellZones_0 0/pointZones_0
    checkMesh -allTopology -allGeometry | tee log.checkMesh
    decomposePar -force                | tee log.decomposePar
fi

# 5. Refresh decomposed initial fields from the (possibly reused) mesh
rm -rf processor*/0 processor*/0.*
decomposePar -force | tee log.decomposePar

# 6. Run the actual solver in parallel
mpirun --use-hwthread-cpus -n 16 rhoReactingBuoyantFoam -parallel | tee log.solver

# 7. Reconstruct only the latest time directory back into the host
reconstructPar -latestTime | tee log.reconstructPar

# 8. Mark this case complete (run_doe.sh checks for SIM_DONE)
echo "[case_04] SIM DONE at $(date +%FT%T)"
touch SIM_DONE
\end{lstlisting}

\subsubsection{The eight things this does}

\begin{enumerate}[leftmargin=1.5em,topsep=2pt]
    \item Source OpenFOAM (sets \texttt{\$WM\_PROJECT\_DIR} etc.).
    \item Read per-case parameters (\texttt{UMAIN, UBRANCH, D2, \dots}).
    \item Mesh-reuse shortcut: if the same-slice donor case has
    already meshed, just \texttt{cp} the mesh (saves $\sim 30$ minutes
    per case).
    \item Otherwise, full mesh build:
    \begin{itemize}[leftmargin=1em,topsep=1pt]
        \item \texttt{generateSTL.py} --- analytic T-junction STL.
        \item \texttt{clip\_stls.py} --- clip to half-domain $x\ge 0$.
        \item \texttt{surfaceFeatureExtract} --- extract sharp edges.
        \item \texttt{blockMesh} --- background mesh.
        \item \texttt{snappyHexMesh -overwrite} --- castellate, snap, layers.
        \item \texttt{createPatch} --- re-tag \texttt{allBoundary} faces.
        \item \texttt{checkMesh} --- mesh-quality report.
        \item \texttt{decomposePar -force} --- split into 16 ranks.
    \end{itemize}
    \item Refresh \texttt{processor*/0} from the source \texttt{0/}
    (ensures field sizes match a possibly reused mesh).
    \item \texttt{mpirun -n 16 rhoReactingBuoyantFoam -parallel} ---
    the actual transient.
    \item \texttt{reconstructPar -latestTime} --- merge the latest
    time directory.
    \item \texttt{touch SIM\_DONE} --- the orchestrator sentinel.
\end{enumerate}

\clearpage
% =====================================================================
\section{The DoE pipeline --- how the 10 cases are generated}
% =====================================================================

\subsection{lhs\_design.py --- the sliced LHS sampler}

\texttt{lhs\_design.py} produces \texttt{doe\_design.csv}, a 10-row
table of $(\dD,\,\HBR,\,\VR)$ triplets:

\begin{lstlisting}[style=pystyle,caption={Sketch of \texttt{lhs\_design.py}}]
def sliced_lhs(seed=42, n_slices=5, m=2):
    rng = np.random.default_rng(seed)
    # 1. uniform LHS over (dD, HBR) box
    dD_levels  = np.linspace(0.15, 0.45, n_slices+1)
    cases = []
    for s, (lo, hi) in enumerate(zip(dD_levels[:-1], dD_levels[1:])):
        # 2. sample 2 HBR values within the slice
        hbr_lhs = lhs(2, m, seed=rng)
        for j in range(m):
            dD  = rng.uniform(lo, hi)
            hbr = 0.05 + 0.15 * hbr_lhs[j]
            # 3. invert HBR(dD, VR) = (dD)^2 * VR / (1 + (dD)^2 * VR)
            VR  = hbr / ((dD**2) * (1 - hbr))
            assert VR <= 12, "VR safety cap exceeded"
            cases.append((s+1, dD, hbr, VR))
    return cases
\end{lstlisting}

\paragraph{The point of the slice structure.}
Each pair of cases shares the same $\dD$ value (slice). Geometry only
depends on $\dD$, so the same mesh works for both. The orchestrator
\texttt{run\_doe.sh} therefore builds a mesh once per slice (5 builds
not 10), saving roughly 4 hours per campaign.

\paragraph{Why the soft cap on VR.}
$\HBR$ is bounded above at 0.20, but for very small $\dD$ the equation
$\HBR = (\dD)^2 \VR / (1+(\dD)^2 \VR)$ gives an enormous $\VR$ to hit
20\% blend. A $\VR>12$ jet impinges on the opposite wall and erodes
the pipe (and triggers very small time-steps in our solver). The
soft cap is a physically motivated pruning of the parameter box.

\subsection{stamp\_cases.py --- materialise the 10 case dirs}

\texttt{stamp\_cases.py} reads \texttt{doe\_design.csv}, copies the
\texttt{doe\_base/} template, and replaces tokens like
\texttt{@UBRANCH@}, \texttt{@D2@}, \texttt{@KBRANCH@}, \texttt{@SA@},
\texttt{@CA@} with the per-case numerics. After stamping, each
\texttt{case\_NN/} directory is a fully populated, ready-to-run
OpenFOAM case (its own \texttt{0/}, \texttt{constant/}, \texttt{system/},
\texttt{Allrun}, \texttt{generateSTL.py}, \texttt{case\_info.json}).

\subsection{run\_doe.sh --- the orchestrator}

The orchestrator iterates over the 10 cases, decides whether to
build a mesh or reuse one, and handles per-case success/failure:

\begin{lstlisting}[style=bashstyle,caption={Core loop of \texttt{run\_doe.sh}}]
for cname in "${SELECTED[@]}"; do
    cdir="${CASES_DIR}/${cname}"
    [[ -d "${cdir}" ]] || continue

    if [[ -f "${cdir}/SIM_DONE" ]]; then
        log "[skip] ${cname}: already done"
        post_process "${cdir}"
        continue
    fi

    # mesh-reuse: find a donor in the same slice
    slice="$(case_slice_id "${cdir}")"
    donor="$(find_mesh_donor "${slice}" "${cdir}")" || true
    [[ -n "${donor}" ]] && export REUSE_MESH_FROM="${donor}" || unset REUSE_MESH_FROM

    if bash "${cdir}/Allrun" >> "${cdir}/log.run_doe" 2>&1; then
        log "[${cname}] SIM OK"
        stamp_wall_seconds "${cdir}/case_info.json" "${elapsed}" ok
        post_process "${cdir}"
    else
        log "[${cname}] SIM FAILED"
        tail -n 200 "${cdir}/log.run_doe" > "${cdir}/SIM_FAILED"
        post_process "${cdir}"
        log "[${cname}] continuing with next case"   # crash resilience
    fi
done
\end{lstlisting}

\paragraph{Crash resilience.}
A failure in one case never aborts the campaign. The orchestrator
copies whatever the case did manage to produce (the partial
\texttt{processor*} state, the log) into \texttt{doe\_results/case\_NN/},
writes a \texttt{SIM\_FAILED} sentinel with the last 200 log lines,
and moves on. This is \emph{exactly} how case~04 of the 30\textdegree{}
campaign was paused on demand and the campaign continued on case~05.

\subsection{generateSTL.py --- the geometry}

We do not draw the T-junction in a CAD package. Instead a Python
script generates a watertight STL parametrically. The cylinders are
sampled on a $(\theta, z)$ grid; the intersection curve between the
branch cylinder and the main cylinder is computed analytically; a
rectangular hole in the main pipe is punched in $(\theta, z)$ space;
and the boundary of that hole is ``zippered'' to the intersection
curve by triangulation. The output is a single STL with $\sim 44\,000$
triangles whose volume agrees with the analytic volume to better
than 1\%.

\paragraph{The 30\textdegree{} bug, in one sentence.}
The original code sized the rectangular hole using
\(z_\mathrm{stretch} = R_2 + 0.5\,L_\mathrm{branch}\,|\!\cos\alpha|\),
which is fine at $\alpha = 90^\circ$ ($\cos\alpha = 0$ collapses the
heuristic), but at $\alpha = 30^\circ$ punches a rectangle
$\sim 7\times$ wider than the actual intersection curve, forcing the
zipper to span $\sim 1$\,m of empty annular space with stretched
triangles that snappyHexMesh interpreted as a wall. The fix is to size
the rectangle directly from the bbox of the intersection curve plus
a few grid cells of buffer.

\clearpage
% =====================================================================
\section{Numerics refresher --- the things you'll be asked}
% =====================================================================

\subsection{Why the time scheme is Euler and not Crank--Nicolson}

\begin{itemize}[leftmargin=1.5em,topsep=2pt]
    \item \textbf{Euler (1st order, implicit):} stable, monotonic,
    cheap. Diffusive in time --- but the time-averaging window
    (0.6--1.2\,s) absorbs the diffusion, so reported QoIs are not
    affected.
    \item \textbf{Crank--Nicolson (2nd order):} less diffusive but
    can produce wiggles in stiff equations (energy, $\omega$). For a
    parametric DoE the bookkeeping cost outweighs the benefit.
    \item \textbf{Backward (2nd order, BDF):} a compromise; we did
    not need it because the time-averaged QoIs converge fine with
    Euler.
\end{itemize}

\subsection{Why upwind and not vanLeer / linearUpwind}

\begin{itemize}[leftmargin=1.5em,topsep=2pt]
    \item \texttt{upwind}: 1st order, monotonic, $O(\Delta x)$
    diffusive.
    \item \texttt{linearUpwind}: 2nd order in smooth regions, can wiggle near
    sharp gradients.
    \item \texttt{vanLeer}: TVD limited 2nd order, bounded but
    noticeably more expensive per step.
\end{itemize}
For this DoE, upwind on $H_2/U/k/\omega$ keeps everything bounded and
cheap; the $\sim 5$--$10\%$ excess CoV from numerical diffusion is
reported transparently (and is the same across all cases, so trends
are unaffected).

\subsection{What CFL / Courant number actually controls}

The local Courant number is
\(\mathrm{Co} = |\mathbf{u}|\,\Delta t / \Delta x\). It measures the
fraction of a cell a fluid parcel crosses in one time step. For
explicit schemes the stability limit is $\mathrm{Co}\le 1$.
For implicit PIMPLE we can run $\mathrm{Co}\sim 5$--$10$ as long as the
PIMPLE inner loop converges per step. Hence
\texttt{maxCo 4.5} (90\textdegree{} campaign) and \texttt{maxCo 8.0}
(30\textdegree{} campaign).

\subsection{What \texttt{p} vs \texttt{p\_rgh} actually buys you}

If we kept solving for $p$ directly, gravity would appear as a body
force in the momentum equation, and the pressure equation would have
to fight against the (large, easily-known) hydrostatic distribution
$\rho \mathbf{g}\!\cdot\!\mathbf{x}$ to find the (small, interesting)
kinematic part. By solving for $p_\mathrm{rgh}=p-\rho\mathbf{g}\!\cdot\!\mathbf{x}$
instead, the hydrostatic component drops out analytically and the
solver only resolves the kinematic pressure --- which is what
matters for our $\Delta p$ calculation. \texttt{p} is reconstructed
post-hoc.

\clearpage
% =====================================================================
\section{Physics refresher --- the dimensionless numbers}
% =====================================================================

\begin{table}[H]
\centering
\small
\caption{The four dimensionless numbers I should know cold.}
\begin{tabular}{l c l l}
\toprule
\textbf{Number} & \textbf{Definition} & \textbf{Meaning} & \textbf{Our regime} \\
\midrule
Reynolds  & $Re = \rho U D / \mu$
          & inertia / viscous
          & $5\times 10^5$ (fully turbulent) \\
Courant   & $\mathrm{Co} = |\mathbf{u}|\Delta t/\Delta x$
          & advection / cell time
          & $\le 4.5$--$8.0$ (numerical stability) \\
Richardson & $\Ri = g\Delta\rho L /(\rho U^2)$
           & buoyancy / shear
           & $\sim 10^{-2}$ (turbulence-dominated) \\
Momentum-flux ratio & $J = \rho_b u_b^2 /(\rho_m u_m^2)$
                    & jet penetration
                    & 0.7--160 across the DoE \\
Froude    & $\Frnum = U/\sqrt{g\Delta\rho L/\rho}$
          & inverse $\sqrt{\Ri}$
          & $\gg 1$ (production), $\sim 0.3$ (sanity) \\
Wobbe Index & $I_W = V_c/\sqrt{G}$
            & energy delivery per unit area
            & not modelled here \\
\bottomrule
\end{tabular}
\end{table}

\subsection{What the project's primary metrics are, mathematically}

\paragraph{Coefficient of Variation (CoV) at outlet.}
\begin{equation}
\CoV \;=\; \frac{\sigma_{Y_{H_2}}}{\overline{Y_{H_2}}}
\;=\; \frac{\sqrt{\sum_i w_i (Y_{H_2,i}-\overline{Y_{H_2}})^2 / \sum_i w_i}}{\sum_i w_i Y_{H_2,i}/\sum_i w_i}.
\end{equation}
Three weightings:
$w_i = A_i$ (area), $w_i = \rho_i\,|\mathbf{u}\!\cdot\!\hat{\mathbf{n}}|A_i$
(mass-flux), $w_i = |\mathbf{u}\!\cdot\!\hat{\mathbf{n}}|A_i$
(volumetric-flux). The mass-flux weighting is the
custody-transfer relevant one (because the customer downstream
receives the mass-weighted average), but all three are reported.
Industry threshold: $\CoV \le 0.05$.

\paragraph{Static pressure drop.}
$|\Delta p| = |p_{\mathrm{main\_inlet}} - p_{\mathrm{outlet}}|$,
patch-area-averaged, time-averaged over the second flow-through.

\paragraph{Danckwerts segregation index.}
\begin{equation}
I_s \;=\; \frac{\overline{(Y_{H_2}-\overline{Y_{H_2}})^2}}{\overline{Y_{H_2}}\cdot(1-\overline{Y_{H_2}})}.
\end{equation}
$I_s = 1$ means completely segregated (no mixing); $I_s = 0$ means
perfectly mixed. Used as a cross-check on CoV because it is bounded.

\clearpage
% =====================================================================
\section{Things I might be asked, with crisp answers}
% =====================================================================

\begin{vivabox}[Q1: Why a T-junction and not a Y-junction or coaxial mixer?]
T-junctions are the dominant industrial topology for hydrogen
injection into existing natural-gas pipelines because they are
\emph{retrofittable}: a hot tap can punch the branch into a live main
without taking it out of service. Y- and coaxial mixers require
custom fabrication and a shutdown. The cost of poorer mixing is what
we are quantifying.
\end{vivabox}

\begin{vivabox}[Q2: Why three angles (90, 30, 150)?]
$90^\circ$ is the industry baseline (cross-flow). $30^\circ$ is
co-flow --- shallowest practical industrial geometry, gentle on the
pipe. $150^\circ$ is counter-flow --- the opposite extreme. Three
points span the regime; the linear/quadratic interpolation across
angle gives you a usable design surface for any $\alpha$ in
$[30^\circ,150^\circ]$.
\end{vivabox}

\begin{vivabox}[Q3: Why sliced LHS instead of full factorial?]
A 3-factor 5-level full factorial is $5^3=125$ cases. We have
budget for 10 per angle. LHS samples \emph{uniformly} on every
1-D projection; sliced LHS additionally pre-stratifies $\dD$ into 5
slices so we get mesh-reuse savings. Statistically, the effect of
each factor is recovered to within $\sim 5\%$ of the full-factorial
sensitivity for low-curvature responses.
\end{vivabox}

\begin{vivabox}[Q4: Why is my DoE only 10 cases? Isn't that small?]
For a quadratic response surface in 3 factors we need 10 distinct
points minimum (Box--Behnken or central composite design with 1
centre point). Sliced LHS at 10 points gives the same coverage with
better space-filling properties. A future objective would be to add
a second wave of 10 cases at the inflection regions identified by
the first wave (sequential DoE), but that is beyond the MTP scope.
\end{vivabox}

\begin{vivabox}[Q5: Symmetry plane --- aren't real T-junctions asymmetric?]
The instantaneous flow is asymmetric (vortex shedding, turbulent
fluctuations), but the \emph{time-averaged} flow is mirror-symmetric
under the $x\!\to\!-x$ reflection because the inlet conditions are.
RANS solves the time-averaged equations directly. Half-domain RANS
of T-junctions has been validated against full-domain RANS in
Chen et al.\ (2017), Frank et al.\ (2010), Walker et al.\ (2010);
the time-averaged QoIs (CoV, $\Delta p$) agree to within $\sim 1\%$.
The symmetry would NOT be valid for LES.
\end{vivabox}

\begin{vivabox}[Q6: Why rhoReactingBuoyantFoam and not rhoPimpleFoam?]
Three things distinguish them:
\begin{enumerate}[leftmargin=1.5em,topsep=2pt]
    \item \texttt{rhoReactingBuoyantFoam} solves species transport
    natively with full multispecies thermo (JANAF, mixture viscosity,
    Sutherland). \texttt{rhoPimpleFoam} only solves enthalpy and a
    single-species $\rho$.
    \item It uses $p_\mathrm{rgh}$ (gravity-corrected pressure), so
    buoyancy is not an explicit body-force term in the momentum
    equation but is built into the pressure decomposition.
    \item It includes a chemistry hook (we set \texttt{combustion off}
    and an empty \texttt{reactions} dictionary, but the solver
    machinery is there if we ever need it).
\end{enumerate}
\end{vivabox}

\begin{vivabox}[Q7: You report CoV at one downstream station. Why not several?]
We \emph{do} sample $H_2$ at four downstream stations
($z = 3, 4, 5, 6\,\si{\meter}$) using \texttt{sampledSurface} function
objects in \texttt{controlDict}. The headline CoV is at the outlet
($z = 6.9\,\si{\meter} = 4.6\,D$ downstream of the junction)
because that's the relevant ``custody-transfer'' station. The
upstream four stations give the \emph{mixing length} --- how far
downstream you'd need to go to drop CoV below a threshold.
\end{vivabox}

\begin{vivabox}[Q8: How do you validate the mesh?]
Three layers:
\begin{enumerate}[leftmargin=1.5em,topsep=2pt]
    \item \texttt{checkMesh -allTopology -allGeometry}: max
    non-orthogonality 49\textdegree{}/65\textdegree{}, max skewness
    2.57/4, layer coverage 99.6\%. All within OpenFOAM healthy
    thresholds.
    \item Grid-independence: 230\,k, 460\,k, 920\,k cells. CoV changed
    by 0.4\% (coarse$\rightarrow$medium) and 0.2\%
    (medium$\rightarrow$fine). Adopted the medium for the campaign.
    \item Single connected fluid region: \texttt{checkMesh} reports
    \texttt{Number of regions: 1 (OK)} and a Python
    \texttt{pyvista.connectivity} check confirms it. (Crucial after
    the 30\textdegree{} mesh-defect fix.)
\end{enumerate}
\end{vivabox}

\begin{vivabox}[Q9: What if the prof asks ``why isothermal?'']
Two reasons:
\begin{enumerate}[leftmargin=1.5em,topsep=2pt]
    \item Industrial gas pipelines run essentially isothermal
    (~5\textdegree{}C variation over kilometres). The Joule--Thomson
    coefficient of $H_2$ is small in our pressure regime and would
    only matter near regulators.
    \item Disabling the energy equation (relTol 1, maxIter 0) saves
    $\sim 20\%$ wall time per case, which over the DoE matters. If we
    later cared about thermal effects we would just re-enable it.
\end{enumerate}
\end{vivabox}

\begin{vivabox}[Q10: What is the limitation of this study?]
Honestly stated:
\begin{itemize}[leftmargin=1.5em,topsep=2pt]
    \item RANS time-averaging hides the unsteady mixing structures
    that LES would resolve.
    \item Wall functions ($y^+\sim 30$--$80$) sacrifice
    near-wall-resolution; a $y^+\sim 1$ run would be the next
    refinement.
    \item No combustion, no chemistry. We are studying mixing only.
    \item Single fluid pair (H$_2$/CH$_4$). A real pipeline has
    multi-component natural gas (CH$_4$, C$_2$H$_6$, N$_2$, CO$_2$).
    \item One operating pressure (6.9\,MPa). High-pressure CFD
    requires real-gas EOS, which we have not used.
    \item Symmetry plane assumption excludes any pitchfork
    instability the geometry may admit at higher $\VR$.
\end{itemize}
\end{vivabox}

\begin{vivabox}[Q11: Why y+ around 30-80 specifically and not just any value?]
The wall-function regime is valid for $y^+ > 11$ (above the
viscous sublayer) and $y^+ < 200$ (still inside the log layer).
$y^+\sim 30$--$80$ keeps us comfortably inside that band on every
mesh and at every cell, with margin for the highest-VR case where
near-wall velocity spikes increase $u_\tau$.
\end{vivabox}

\begin{vivabox}[Q12: How do you know the simulation is converged?]
Three checks per case:
\begin{itemize}[leftmargin=1.5em,topsep=2pt]
    \item Final residuals: $\mathbf{u}\sim 10^{-9}$,
    $p_\mathrm{rgh}\sim 10^{-7}$, $H_2\sim 10^{-8}$.
    \item Cumulative continuity error $\sim 5\times 10^{-3}$
    ($<10^{-2}$ acceptance).
    \item Time-traces of outlet $H_2$, $\Delta p$ flatten over the
    second flow-through (the moving average is constant to $\sim 1\%$
    over the averaging window).
\end{itemize}
\end{vivabox}

\begin{vivabox}[Q13: Why 1.2 s endTime?]
Two main-pipe flow-throughs ($L_\mathrm{main}/U_\mathrm{main} =
6.9/10 = 0.69\,\si{\second}$). The first one is the start-up transient
(filling the pipe with the inlet condition), the second is the
window we time-average over. This is standard practice for transient
RANS of internal pipe flow (Ayach 2017).
\end{vivabox}

\begin{vivabox}[Q14: What does ``mesh-reuse from case\_03'' mean?]
A sliced LHS pairs cases by $\dD$. Cases 03 and 04 both have
$\dD = 0.252$, so they have the \emph{same} geometry and \emph{same}
mesh. Building the mesh takes $\sim 30$ minutes (\texttt{snappyHexMesh}
+ layers); reusing it means just \texttt{cp -r constant/polyMesh} and
the \texttt{processor*/constant/polyMesh} sub-meshes from the donor.
This trick saves about 4 hours over the 10-case campaign.
\end{vivabox}

\begin{vivabox}[Q15: What happens if a case crashes mid-run?]
The hardened \texttt{Allrun} (in spirit) and the orchestrator's
failure handler ensure:
\begin{itemize}[leftmargin=1em,topsep=2pt]
    \item \texttt{reconstructPar -latestTime} runs whatever was
    written before the crash;
    \item \texttt{SIM\_FAILED} is written with the last 200 log lines;
    \item the campaign moves on to the next case.
\end{itemize}
This is exactly what happened on the 30\textdegree{} case~04 when we
manually paused it to free the queue. The partial state at $t=0.2$\,s
is still on disk and can be resumed via the dedicated
\texttt{resume\_case04\_30deg.sh} script.
\end{vivabox}

\begin{vivabox}[Q16: Why is the operating pressure 6.9 MPa?]
That is the pipeline operating pressure cited by Su et al.\ (2022)
for a typical natural-gas trunk line. We match it so our results
sit in the same operating point as the literature we benchmark
against.
\end{vivabox}

\begin{vivabox}[Q17: Why stop the campaign at $\alpha = 30^\circ$ and $150^\circ$
and not $0^\circ$ or $180^\circ$?]
$\alpha = 0^\circ$ would mean injecting along the main pipe axis from
upstream --- physically just adding the gases through the same inlet,
with no T-junction. $\alpha = 180^\circ$ would mean injecting from
downstream backward against the flow, also unphysical for an
injection problem. The practical industrial range is $\sim
30^\circ$--$150^\circ$ depending on terrain and ROW constraints.
\end{vivabox}

\clearpage
% =====================================================================
\section{Quick command-line cheatsheet}
% =====================================================================

\begin{lstlisting}[style=bashstyle,caption={The 25 commands that cover 95\% of OpenFOAM use}]
# 1. Source the environment (do this in every fresh shell)
source /usr/lib/openfoam/openfoam2406/etc/bashrc

# 2. Inspect a case
foamGet -case <case>           # generate boilerplate dictionary
checkMesh -allTopology         # mesh quality
checkMesh -allGeometry
foamListTimes -case <case>     # which time directories exist
postProcess -func sample       # sample fields on planes/lines

# 3. Build the mesh
blockMesh                       # background mesh from blockMeshDict
surfaceFeatureExtract           # extract eMesh feature edges from STLs
snappyHexMesh -overwrite        # cut + snap + layers in one go
checkMesh                       # always check after meshing
createPatch -overwrite          # rename / merge patches

# 4. Decompose / reconstruct
decomposePar -force             # split for parallel
reconstructPar -latestTime      # merge latest time
reconstructPar -time 0.5:1.2    # merge a range

# 5. Run the solver
mpirun -n 16 rhoReactingBuoyantFoam -parallel
foamLog log.solver              # extract residual time series

# 6. Post-process
paraFoam -touch                 # create the .foam stub for ParaView
postProcess -func 'patchAverage(name=outlet, fields=(H2 p_rgh))'

# 7. Manage the case
foamCleanCase                   # remove time dirs and processor*
wclean libso ./mySolver         # clean a custom solver
wmake libso ./mySolver          # compile a custom solver

# 8. The DoE pipeline (this project specifically)
python3 lhs_design.py --seed 42 --outdir doe/
python3 stamp_cases.py --design doe/doe_design.csv --base . --cases ../cases
./run_doe.sh                    # the orchestrator, in a tmux session
./tools/doe_progress_30deg.sh   # snapshot of remote campaign
\end{lstlisting}

\clearpage
% =====================================================================
\section{Final pre-viva checklist}
% =====================================================================

\textbf{Things I should have ready to draw on a whiteboard.}

\begin{enumerate}[leftmargin=1.5em]
    \item The T-junction geometry (main pipe horizontal, branch coming
    in at $\alpha$, label $D$, $d$, $\dD$, $\VR$, $\HBR$, the symmetry
    plane).
    \item The 1/7-th power-law profile shape.
    \item A jet-in-crossflow sketch with the counter-rotating vortex
    pair, label penetration depth $\propto \sqrt{J}\,d$.
    \item The Richardson / momentum-flux number axes, with
    ``stratified vs.\ mixed'' bands marked.
    \item The PIMPLE loop diagram (predictor $\rightarrow$ pressure
    correction $\times$ nCorrectors $\rightarrow$ outer corrector
    $\times$ nOuterCorrectors).
    \item The case directory tree (\texttt{0/}, \texttt{constant/},
    \texttt{system/} with one example file each).
\end{enumerate}

\textbf{Things I should know by heart.}

\begin{itemize}[leftmargin=1.5em]
    \item Solver = \texttt{rhoReactingBuoyantFoam}; turbulence =
    $k$-$\omega$ SST; thermo = JANAF + Sutherland + perfectGas.
    \item Domain: $D = 0.46$\,m, $L_\mathrm{main} = 6.9$\,m,
    $\dD\in[0.15,0.45]$, half-domain symmetry ($\sim 460$\,k cells).
    \item Time: $t_\mathrm{end} = 1.2$\,s, $\mathrm{Co}_\mathrm{max} = 4.5$
    (90\textdegree{}) or 8.0 (30\textdegree{}), $\Delta t \sim 10$\,$\mu$s.
    \item DoE: sliced LHS, seed 42, 10 cases per angle (5 slices of 2),
    same parameter set across angles for clean factorial comparison.
    \item Three angles: 90\textdegree{} (done), 30\textdegree{}
    (campaign in progress, mesh-fix story), 150\textdegree{} (planned).
    \item 90\textdegree{} headline result: best CoV = 0.066 at
    $\VR=3.81$ (case~04); worst CoV = 0.71 at $\VR=0.69$ (case~08).
\end{itemize}

\textbf{Things to NOT bluff if asked.}

\begin{itemize}[leftmargin=1.5em]
    \item I have not done LES. RANS only. If asked about LES, say I
    expect it to break the symmetry assumption and that we'd need
    full-domain meshes.
    \item I have not validated against experiment, only against
    Su 2022 trends (which is qualitative, not absolute-number, agreement).
    \item The 30\textdegree{} and 150\textdegree{} campaigns are at
    different stages of completion; the cross-angle synthesis is the
    final-viva deliverable, not done yet.
    \item Real gas effects (Lee--Kesler, Peng--Robinson) are not
    modelled. Operating pressure is well below the H$_2$ critical
    point, so this is an acceptable approximation but not physics-tight.
\end{itemize}

\end{document}
